% Generated draft: requires mathematical and source review.
\documentclass[11pt]{article}
\usepackage[margin=25mm]{geometry}
\usepackage{fontspec}
\setmainfont{Latin Modern Roman}
\usepackage{amsmath,amssymb,mathtools,unicode-math}
\setmathfont{Latin Modern Math}
\usepackage{enumitem,needspace,xurl,hyperref}
\setlength{\parindent}{0pt}
\setlength{\parskip}{5pt}
\setlength{\emergencystretch}{3em}
\newcommand{\N}{\mathbb{N}}
\newcommand{\Z}{\mathbb{Z}}
\newcommand{\Q}{\mathbb{Q}}
\newcommand{\R}{\mathbb{R}}
\newcommand{\C}{\mathbb{C}}
\newcommand{\F}{\mathbb{F}}
\newcommand{\A}{\mathbb{A}}
\newcommand{\PP}{\mathbb{P}}
\newcommand{\E}{\mathbb{E}}
\newcommand{\eps}{\varepsilon}
\newcommand{\dd}{\,\mathrm{d}}
\begin{document}
\section*{UnsolvedMath ID: 7200068}
\textbf{Source problem code:} catalogue entry
\subsection{Definitions and mathematical statement}
A holyhedron is conventionally a polyhedron every face of which has at least one hole. Here ``face'' must permit a planar polygonal region with nontrivial holes, rather than requiring every face to be a topological disk. For a specified admissible holyhedron \(P\), let \(\mathcal F(P)\) be its face set. Determine
\[
 f_{\min}=\min\{|\mathcal F(P)|:P\text{ is an admissible holyhedron}\}.
\]
A determination requires both a construction achieving the claimed number and a proof that every admissible holyhedron has at least that many faces.

The supplied excerpt does not define holyhedron or its polyhedral conventions. The usual meaning of a hole in a planar face is a bounded complementary component within its supporting plane, but embeddedness, closedness, connectedness, manifold requirements, and whether faces are maximal connected planar patches must be fixed. Arbitrarily subdividing faces or grouping coplanar pieces could otherwise change the count. No admissible class is silently selected beyond the conventional description above.

The summary reports finite-face examples exist, attributes existence to J. H. Conway, and says the exact minimum is unknown, without giving any construction, face count, or authoritative citation. The existence attribution is not certified here. Its report of no 2024--2026 minimum determination and the pinned Wikipedia open label remain source assessments pending primary-source review.
\subsection{Short English statement}
What is the smallest face count of a polyhedron whose every face has a hole? The face and polyhedron conventions need confirmation, and the record supplies neither a numerical bound nor an authoritative construction reference.
\subsection{Sources}
\begin{itemize}
\item Local registry: \url{data/unsolvedmath/problems.json}, permanent UnsolvedMath ID \(7200068\), source code catalogue entry; supplied statement, background, and research summary, updated 2026-08-06.
\item Wikipedia contributors, \emph{List of unsolved problems in mathematics}, Geometry bullet 68, pinned revision 1366636767: \url{https://en.wikipedia.org/w/index.php?title=List_of_unsolved_problems_in_mathematics&oldid=1366636767#Geometry}. MediaWiki-source extraction accessed 2026-07-29; primary-source status review remains required.
\item J. H. Conway is named by the supplied summary for existence, without an original publication, construction, or face count. No authoritative minimum-bound reference is supplied.
\end{itemize}
\end{document}
